\section{Related literature}\label{sec:literature}

\paragraph{Economics of cybercrime.} The view that information-security outcomes are shaped by incentives rather than by technology alone goes back to \citet{anderson2001} and \citet{anderson2006}; \citet{moore2009} survey the economics of online crime, emphasizing the specialization of criminal supply chains and the role of cash-out frictions. Our attacker model is a direct application of the rational-choice approach to crime of \citet{becker1968}: the attacker weighs the expected gain against the probability and severity of punishment. Our addition is to separate the expected gain into a breach component and a \emph{retention} component, and to show that the retention component, governed by the irreversibility of settlement and by laundering frictions, dominates the comparison between traditional and crypto-asset venues. \citet{aldasoro2022} and \citet{jamilov2021} document cross-sectional drivers and the economic anatomy of cyber risk with incident and textual data; our model is complementary in providing a structural account of venue choice.

\paragraph{Cyber loss distributions and operational risk.} A consistent empirical finding is that cyber losses are heavy-tailed. \citet{maillart2010} report power-law tails for personal-data breaches, \citet{edwards2016} examine breach sizes and frequencies and caution against reading trends into heavy-tailed data, and \citet{eling2019} and \citet{biener2015} fit extreme-value models to operational-loss databases and find that cyber events are distinctive mainly in their tails. Cyber risk in banks has historically been modelled within operational risk, where the loss distribution approach, which combines a frequency and a severity distribution into an aggregate loss, was standard under the advanced measurement approaches \citep{frachot2001,moscadelli2004,chavez2006}, before being replaced by a standardized approach in the final Basel III framework \citep{bcbs2017}. \citet{moscadelli2004} and \citet{chavez2006} already noted that tail estimates from operational-loss data are fragile. We quantify that fragility for interval estimation specifically and show that it is one-sided.

\paragraph{Extreme-value inference.} The peaks-over-threshold method rests on the Pickands--Balkema--de Haan theorem \citep{balkema1974,pickands1975}: for a wide class of distributions, exceedances of a high threshold are approximately generalized Pareto. Maximum-likelihood estimation and its asymptotics are due to \citet{smith1987} and \citet{davison1990}; \citet{hosking1987} propose probability-weighted moments and \citet{zhang2009} an efficient Bayesian-flavoured alternative; \citet{grimshaw1993} reduces the likelihood to a one-dimensional profile, which we use. The Hill estimator \citep{hill1975} provides a non-parametric diagnostic of the tail index. Textbook treatments \citep{embrechts1997,coles2001,mcneil2015} note that likelihood-ratio intervals for return levels are usually preferable to Wald intervals because the profile likelihood is strongly asymmetric. To our knowledge, the magnitude and the one-sidedness of the coverage failure at record lengths typical of cyber data, and the comparison with the parametric bootstrap \citep{davison1997} under a proper scoring rule \citep{gneiting2007}, have not been documented for cyber or operational losses. We also contribute a simple second-order correction for compound sums, building on the single-big-jump structure of subexponential sums \citep{embrechts1997}, and we use Breiman's lemma \citep{breiman1965} to establish the tail index of our loss process.

\paragraph{Financial networks and systemic cyber risk.} The clearing model of \citet{eisenberg2001} and its successors characterize how defaults propagate through obligations; \citet{acemoglu2015} show that the same network can be robust to small shocks and fragile to large ones, and \citet{glasserman2015,glasserman2016} bound the amplification that balance-sheet contagion can generate. DebtRank \citep{battiston2012} and its linear variant \citep{bardoscia2015} propagate distress before default by marking down claims in proportion to the debtor's equity loss; \citet{bardoscia2017} show that such systems become unstable when the spectral radius of the exposure matrix exceeds one. For cyber risk specifically, \citet{eisenbach2022} estimate how an attack that impairs a large bank's payment capacity would spill over to other banks through the wholesale payment network, \citet{duffie2019} analyse the liquidity runs a cyber event could trigger, \citet{kashyap2019} propose principles for regulating cyber risk, and \citet{crosignani2023} show that attacks propagate through supply chains. Our network analysis contributes a shared-vendor scenario, the third-party concentration channel targeted by recent regulation \citep{dora2022}, and a demonstration that interconnection, not size, drives systemic importance in a core--periphery system.
